\documentclass[oneside]{scrbook}
% \documentclass{article}

% \usepackage[utf8]{inputenc}
\usepackage[T1]{fontenc}
\usepackage{version}
% \usepackage[latin1]{inputenc}
\usepackage[italian]{babel}
\usepackage{algpseudocode}
\usepackage{amsmath}
\usepackage{booktabs}
\usepackage{color}
\usepackage{amsfonts}
\usepackage{url}
\usepackage{algorithm}
\usepackage{chngcntr}
\usepackage[answerdelayed]{exercise}
\usepackage{graphicx}
\usepackage{caption}
\usepackage{listings}
\usepackage{tikz}
\usepackage{pgf}
\usepackage{colortbl}
\usepackage{listingsutf8}
\usepackage[pdftitle = {Dispense di esercizi}, pdfauthor = {Gianmarco Tedeschi e Alberto Marchesi}, pdfsubject = {Fondamenti di Informatica ($082746$ - A.A.~$2025-2026$)}, pdfcreator = {Gianmarco Tedeschi e Alberto Marchesi}, pdfproducer = {-}]{hyperref}
% \usepackage[automark,nouppercase]{scrpage2}
\usepackage{mdframed}
\usepackage{scrlayer-scrpage} % Necessary package for page styles

%\pagestyle{scrheadings}

%\usepackage{titlesec}
%\newcommand{\sectionbreak}{\clearpage}
\counterwithin{Exercise}{chapter}
\counterwithin{Answer}{chapter}

\algrenewcommand\algorithmicthen{\textbf{Allora}}
\algrenewcommand\algorithmicif{\textbf{Se}}
\algrenewcommand\algorithmicelse{\textbf{Altrimenti}}
\algrenewcommand\algorithmicend{\textbf{Chiudi}}
\algrenewcommand\algorithmicwhile{\textbf{Finch\'e}}
\algrenewcommand\algorithmicdo{\textbf{Esegui}}
\algrenewcommand{\algorithmiccomment}[1]{\hskip3em\texttt{//#1}}

\setlength{\parindent}{0in}
\setlength{\parskip}{10pt}

\newmdenv[%
        outerlinewidth=2,%
        roundcorner=10pt,%
        leftmargin=0,%
        rightmargin=0,%
        outerlinecolor=black,
        innertopmargin=\topskip,%
        splittopskip=\topskip,%
]{algo}


\ohead{\headmark}
\ihead{}
\chead{}
\ifoot{A.A.~$2025-2026$}
\cfoot{Fondamenti di Informatica (082746) - G. Tedeschi}
\ofoot{Pag.~\pagemark}
\setheadsepline{.4pt}
\setfootsepline{.4pt}

\definecolor{lbcolor}{rgb}{0.9,0.9,0.9}

\lstdefinestyle{cstyle}{%
  inputencoding=utf8/latin1,
  breaklines=true,
  breakatwhitespace=false,
  basicstyle=\small\ttfamily,
  language=c,
  frame=single,
  showstringspaces=false
}

\lstdefinestyle{cstylesmall}{%
  inputencoding=utf8/latin1,
  breaklines=true,
  breakatwhitespace=false,
  basicstyle=\scriptsize\ttfamily,
  language=c,
  frame=single,
  showstringspaces=false
}

\lstdefinestyle{console}{%
  inputencoding=utf8/latin1,
  breaklines=true,
  breakatwhitespace=false,
  basicstyle=\scriptsize\ttfamily,
  showstringspaces=false,
  backgroundcolor=\color{lbcolor}
}


\lstnewenvironment{clang}{\lstset{style=cstyle}}{}
\lstnewenvironment{clangsmall}{\lstset{style=cstylesmall}}{}

\addtocounter{chapter}{0}

\begin{document}

\chapter{Rappresentazione dell'Informazione e Logica Booleana}


\section{Rappresentazione dell'Informazione}

\subsection{Sistemi di numerazione}
Un numero può essere rappresentato in diverse \textbf{basi numeriche}.  
La base $b$ determina il numero di simboli disponibili per rappresentare le cifre.  

\begin{itemize}
  \item \textbf{Base 2 (binaria)}: cifre $\{0,1\}$.
  \item \textbf{Base 8 (ottale)}: cifre $\{0,1,\dots,7\}$.
  \item \textbf{Base 10 (decimale)}: cifre $\{0,1,\dots,9\}$.
  \item \textbf{Base 16 (esadecimale)}: cifre $\{0,1,\dots,9,A,B,C,D,E,F\}$.
\end{itemize}

Un numero $N$ in base $b$ si scrive come:
\[
N = a_{n-1}b^{n-1} + a_{n-2}b^{n-2} + ... + a_1 b + a_0
\]
dove $a_i$ è la cifra nella posizione $i$ ($0 \leq a_i < b$).

\textbf{Esempio:}
\[
(1011)_2 = 1\cdot 2^3 + 0\cdot 2^2 + 1\cdot 2^1 + 1\cdot 2^0 = (11)_{10}
\]

\paragraph{Quante cifre servono in base $b$?}  
Con $n$ cifre in base $b$ si possono rappresentare $b^n$ combinazioni diverse.  
Se vogliamo rappresentare $N$ informazioni distinte, il numero minimo di cifre richiesto è:
\[
n \geq \lceil \log_b(N) \rceil
\]

\textbf{Esempi:}
\begin{itemize}
  \item Per rappresentare $100$ simboli diversi in base 2:  
  $n \geq \lceil \log_2(100) \rceil = \lceil 6.64 \rceil = 7$.  
  Servono quindi almeno 7 bit.
  \item Per rappresentare $500$ simboli diversi in base 16:  
  $n \geq \lceil \log_{16}(500) \rceil = \lceil 2.7 \rceil = 3$.  
  Servono quindi 3 cifre esadecimali.
\end{itemize}

\subsection{Numeri Naturali}

\subsubsection{Conversioni tra basi}
\paragraph{Da decimale a base $b$}
Si usa il metodo delle \textbf{divisioni successive}.
\begin{enumerate}
    \item Si divide il numero $N$ per la base $b$.
    \item Si prosegue con le divisioni finchè non arrivo a quoziente $0$.
    \item Si raccolgono i resti ($r$) dalle divisioni successive in ordine inverso rispetto a come sono stati ottenuti.
    \item La sequenza di resti in base $b$ rappresenterà il numero $N$ nella codifica cercata.
\end{enumerate}

\textbf{Esempio:}  
Convertire $45_{10}$ in binario:

\begin{center}
\begin{tabular}{r|r}
    \multicolumn{2}{c}{2} \\
    45 & 1 \\
    22 & 0 \\
    11 & 1 \\
    5 & 1 \\
    2 & 0 \\
    1 & 1 \\
    0 &
\end{tabular}
\end{center}

I resti presi in ordine inverso rispetto a quello con cui vengono ottenuti danno come risultato: $45_{10} = 101101_2$.


\paragraph{Da base $b$ a decimale}
Si usa il metodo della \textbf{somma di prodotti}. Se un numero in base $b$ è
$$
    (a_n a_{n-1}...a_1 a_0)_b
$$
dove ogni cifra $a_i$ è compresa tra 0 e $b-1$, allora il valore decimale si calcola come
$$
    N_{10} = a_n \cdot b^n + a_{n-1} \cdot b^{n-1} +...+a_1 \cdot b^1 +a_0 \cdot b^0.
$$
Si andrà quindi a moltiplicare ogni cifra per una potenza della base corrispondente alla sua posizione sommando infine tutti i risultati ottenuti.

\textbf{Esempio:}  
$(253)_8 = 2\cdot 8^2 + 5\cdot 8^1 + 3\cdot 8^0 = 171_{10}$

\paragraph{Conversioni tra basi potenze di 2}
Entrambi i procedimenti esposti hanno validità generale e permettono la conversione da una qualsiasi base $b_i$ a $b_j$. Solitamente, il procedimento delle divisioni successive si utilizza per il passaggio da base 10 a una qualsiasi base $b$. Differentemente, il procedimento della somma di prodotti risulta immediato per la conversione da una qualsiasi base $b$ a base 10. Nel caso di conversioni in cui la base 10 non sia coinvolta è consuetudine effettuare un doppio passaggio, da una base generica a base 10 e da base 10 alla base desiderata. Questo passaggio intermedio attraverso la base 10 può essere omesso, quando si effettua una conversione di base tra due basi tali che una è una potenza dell’altra. Passare dalla base 2 (binario) alla base 8 (ottale) o 16 (esadecimale) è particolarmente semplice, perché $8 = 2^3$ e $16 = 2^4$.  
Questo significa che ogni cifra ottale corrisponde esattamente a un gruppo di 3 bit, e ogni cifra esadecimale corrisponde a un gruppo di 4 bit.

\begin{enumerate}
  \item Si parte dalla cifra meno significativa (a destra).
  \item Si raggruppano i bit a blocchi di 3 (per l’ottale) o 4 (per l’esadecimale).
  \item Se l’ultimo gruppo (quello più a sinistra) non è completo, si aggiungono zeri all’inizio per completarlo.
  \item Si converte ciascun gruppo nella cifra corrispondente.
\end{enumerate}

\textbf{Esempio 1: da binario a ottale}  
Convertire $(101101)_2$ in base 8.  
\[
(101101)_2 = (101)(101)_2
\]  
Ogni gruppo di 3 bit corrisponde a una cifra ottale:  
$101_2 = 5_8$, $101_2 = 5_8$.  
\[
(101101)_2 = (55)_8
\]

\textbf{Esempio 2: da ottale a binario}  
Convertire $(253)_8$ in base 2.  
\[
2 = 010_2, \quad 5 = 101_2, \quad 3 = 011_2
\]  
\[
(253)_8 = (010101011)_2
\]

\subsection{Aritmetica Binaria}
\subsubsection{Somma}
L'operazione di somma è determinata da due elementi principali: il bit di \textbf{risultato} ed il bit di \textbf{resto} definiti come segue:
\begin{center}
    \begin{minipage}[t]{0.45\textwidth}
        \centering
        \textbf{risultato}
        \[
        \begin{array}{c|cc}
             + & 0 & 1 \\
            \hline
            0 & 0 & 1 \\
            1 & 1 & 0
        \end{array}
        \]
    \end{minipage}%
    %\hspace{0.1\textwidth}% Spazio tra le due tabelle
    \begin{minipage}[t]{0.45\textwidth}
        \centering
        \textbf{riporto}
        \[
        \begin{array}{c|cc}
             + & 0 & 1 \\
            \hline
            0 & 0 & 0 \\
            1 & 0 & 1
        \end{array}
        \]
    \end{minipage}
\end{center}
Genericamente, dati due operandi di $n$ bit, $A = [a_{n-1}...a_1a_0]$ e $B = [b_{n-1}...b_1b_0]$ definiamo la somma come:

\[
\begin{array}{*7{c}c} 
    c_{\,n} & c_{\,n-1} & \cdots & c_{\,3} & c_{\,2} & c_{\,1} & \phantom{a_0} & \\
    \phantom{c_n} & a_{\,n-1} & \cdots & a_{\,3} & a_{\,2} & a_{\,1} & a_{\,0} & + \\
    & b_{\,n-1} & \cdots & b_{\,3} & b_{\,2} & b_{\,1} & b_{\,0} & = \\
    \cline{2-7}
    \phantom{+} & s_{\,n-1} & \cdots & s_{\,3} & s_{\,2} & s_{\,1} & s_{\,0} &
\end{array}
\]

Si noti come, nel caso in cui il bit di riporto $c_n = 1$, non è possibile rappresentare il risultato in $n$ bit. Questo fenomeno viene definito come \textbf{overflow}.

\paragraph{Overflow}
Due casi possibili dopo somma:
\begin{enumerate}
    \item Il risultato è rappresentabile con il numero di bit a disposizione, $c_n = 0$.
    \item Il risultato non è rappresentabile con il numero di bit a disposizione, $c_n = 1$.
\end{enumerate}

\subsubsection{Sottrazione}
L'operazione di sottrazione è determinata da due elementi principali: il bit di \textbf{risultato} ed il bit di \textbf{prestito} definiti come segue:
\begin{center}
    \begin{minipage}[t]{0.45\textwidth}
        \centering
        \textbf{risultato}
        \[
        \begin{array}{c|cc}
             - & 0 & 1 \\
            \hline
            0 & 0 & 1 \\
            1 & 1 & 0
        \end{array}
        \]
    \end{minipage}%
    %\hspace{0.1\textwidth}% Spazio tra le due tabelle
    \begin{minipage}[t]{0.45\textwidth}
        \centering
        \textbf{prestito}
        \[
        \begin{array}{c|cc}
             - & 0 & 1 \\
            \hline
            0 & 0 & 0 \\
            1 & 1 & 0
        \end{array}
        \]
    \end{minipage}
\end{center}
Genericamente, dati due operandi di $n$ bit, $A = [a_{n-1}...a_1a_0]$ e $B = [b_{n-1}...b_1b_0]$ definiamo la sottrazione come:

\[
\begin{array}{*7{c}c} 
    p_{\,n} & p_{\,n-1} & \cdots & p_{\,3} & p_{\,2} & p_{\,1} & \phantom{a_0} & \\
    \phantom{c_n} & a_{\,n-1} & \cdots & a_{\,3} & a_{\,2} & a_{\,1} & a_{\,0} & - \\
    & b_{\,n-1} & \cdots & b_{\,3} & b_{\,2} & b_{\,1} & b_{\,0} & = \\
    \cline{2-7}
    \phantom{+} & s_{\,n-1} & \cdots & s_{\,3} & s_{\,2} & s_{\,1} & s_{\,0} &
\end{array}
\]

Analogamente alla somma, in caso il bit di prestito $p_n = 1$ si verifica una situazione anomala definita come \textbf{underflow}.

\subsection{Numeri Relativi}
I numeri relativi (in base 2) possono essere codificati utlizzando diverse rappresentazioni: 
\begin{enumerate}
    \item Rappresentazione in modulo e segno.
    \item Notazione in complemento a 2
\end{enumerate}

\subsubsection{Modulo e Segno}
La rappresentazione in modulo e segno si basa sull'utilizzo di due elementi, \textbf{segno} e \textbf{valore}.
In base 2 si utilizza:
\begin{itemize}
    \item 0 per indicare il segno positivo.
    \item 1 per indicare in segno negativo. 
\end{itemize}
Il bit di segno occupa la posizione più significativa. Ovvero, dato un elemento $N = a_{n-1}a_{n-2}...a_1a_0$, il bit $a_{n-1}$ rappresenterà il segno, i restanti il valore assoluto. 

\textbf{Esempio:}  
Convertire $-14_{10}$ in binario modulo e segno:

\begin{center}
\begin{tabular}{r|r}
    \multicolumn{2}{c}{2} \\
    14 & 0 \\
    7 & 1 \\
    3 & 1 \\
    1 & 1 \\
    0
\end{tabular}
\end{center}

Ricordandoci di aggiungere il bit di segno il risultato sarà: $-14_{10} = 11110_{2\text{MS}}$

Utilizzando questo tipo di rappresentazione per convetire un numero $N$ avremo bisogno di $n = \lceil \log_b(|N| + 1) +1 \rceil$ bit. Ponendoci in binario (base 2) l'insieme dei valori rappresentabili è $[-(2^{n-1}-1), 2^{n-1}-1]$.

\subsubsection{Complemento a 2}
Il complemento a 2 consente di rappresentare numeri negativi in modo compatto, evitando di fornire due rappresentazioni diverse per lo zero e di conseguenza aumentando anche anche i valori rappresentabili. Con $n$ bit (in base 2) si possono infatti rappresentare i numeri nell’intervallo:
\[
[-2^{n-1}, \; 2^{n-1}-1]
\]

I domini di rappresentabilità possono essere definiti come:
\begin{itemize}
    \item Se $N \geq 0$, $n = \lceil \log_b(|N| + 1) +1 \rceil$
    \item Se $N < 0$, $n = \lceil \log_b(|N|) +1 \rceil$
\end{itemize}

Per i numeri positivi si utilizza la stessa metodologia di conversione utilizzata per i numeri naturali, adattata al numero di bit corretto. 
Per i numeri negativi si effettuano i seguenti passaggi:
\begin{enumerate}
  \item Scrivere il numero, in valore assoluto, in binario con $n$ bit.
  \item Invertire tutti i bit (complemento a 1).
  \item Aggiungere $1$.
\end{enumerate}

\textbf{Esempio:}  
Convertire $-12_{10}$ in binario complemento a 2:
Per rappresentare il valore serviranno $n = \lceil \log_2(|-12|) +1 \rceil = 5$ bit.
\begin{center}
\begin{tabular}{r|r}
    \multicolumn{2}{c}{2} \\
    12 & 0 \\
    6 & 0 \\
    3 & 1 \\
    1 & 1 \\
    0
\end{tabular}
\end{center}
$|N| = 12_{10} = 01100_{2\text{C2}}$
\begin{enumerate}
    \item Inverto tutti i bit  $\xrightarrow[]{}10011$ (complemento a 1)
    \item Sommo 1 $\xrightarrow[]{}10100$ (complemento a 2)
\end{enumerate}
Il risultato finale è $N = 10100_{2\text{C2}}$

\paragraph{Somma e Overflow}
Le somme (e sottrazioni) in complemento a 2 si eseguono normalmente in binario.  
In alcuni casi, l’eventuale riporto non comporta una condizione di overflow e può essere ignorato mantenendo il risultato finale corretto.  

\textbf{Esempio:}  
$5 + (-5)$ su 8 bit:  
\[
00000101 + 11111011 = 1\;00000000
\]
Ignoro il riporto: $00000000 = 0$ 

In generale, dati due operandi $a$ e $b$ distinguiamo:
\begin{itemize}
    \item $a>0$, $b>0$ $a+b< 0 \xrightarrow[]{}$ Overflow!
    \item $a>0$, $b<0$ $a+b \xrightarrow[]{}$ Possiamo scartare il bit di carry (se necessario).
    \item $a<0$, $b<0$ $a+b> 0 \xrightarrow[]{}$ Overflow! 
\end{itemize}


\subsection{Numeri Razionali}

\subsubsection{Virgola fissa}
La rappresentazione a virgola fissa riserva:
\begin{itemize}
  \item $n$ bit per la parte intera,
  \item $k$ bit per la parte frazionaria.
\end{itemize}

Per effettuare la conversione di un numero frazionario in binario:
\begin{itemize}
    \item La parte intera si ottiene con il metodo delle divisioni successive.
    \item La parte frazionaria si ottiene con il metodo delle \textbf{moltiplicazioni successive}.
\end{itemize}
Si moltiplica la parte frazionaria per 2 prendendo la parte intera (che sarà 0 o 1). Si rieffettua la stessa operazione sulla nuova parte frazionaria ottenuta dalla moltiplicazione. Il processo si ripete finchè non si raggiunge il numero massimo di bit riservati per la parte frazionaria o quando la parte frazionaria è 0.

\textbf{Esempio: } Convertire $0.375_{10}$ in base 2 utilizzando la virgola fissa.

$0.375 \cdot 2 = 0.75 \to 0$ \\ 
$0.75 \cdot 2 = 1.5 \to 1$  \\
$0.5 \cdot 2 = 1.0 \to 1$  \\

Quindi $0.375 = 0.011_2$.


\subsubsection{Virgola mobile (IEEE 754)}
La virgola mobile ci permette di "spostare" la virgola a seconda del valore che vogliamo codificare.
Il processo si divide in tre fasi: determinazione del \textbf{segno}, calcolo dell'\textbf{esponente} traslato e calcolo della \textbf{mantissa}.

Utilizziamo come esempio il numero $N = +23.75$

\textbf{Fase del segno (S)}: il primo bit è dedicato al segno, 0 se positivo 1 se negativo. Quindi, $S = 0$.

\textbf{Fase di normalizzazione e mantissa (M)}: Il numero decimale (ignorando il segno) deve essere convertito in binario e poi scritto in notazione scientifica binaria normalizzata.
\begin{itemize}
    \item Si converte il numero decimale in binario come per la virgola fissa: 
    
    Parte intera $23_{10} = 10111$

    Parte frazionaria $0.75$
    
    $0.75 \cdot 2 = 1.5 \to 1$  \\
    $0.5 \cdot 2 = 1.0 \to 1$  \\
    
    $23.75_{10} = 10111.11_2$.
    \item Normalizzazione: si sposta la virgola binaria fino a quando non si trova dopo il primo 1. Il numero di posizione in cui la virgola è stata spostata definisce l'\textbf{esponente effettivo}. Spostare la virgola significa moltiplicare piuttosto che dividere per una potenza di 2.

    $E_{eff} = 4$.

    Otteniamo così la mantissa implicita $1.011111$
    \item La mantissa $M$ memorizzata è solo la parte frazionaria che segue il bit implicito (bit prima della virgola).

    $M = 01111100000000000000000$ (23 bit)
\end{itemize}

\textbf{Fase dell'esponente traslato (E)}: Dato l'esponente effettivo possiamo ricavare l'esponente come:
$$
    E = E_{eff} + \text{Bias}
$$
$$
    E = 4 + 127 = 131
$$
Si converte infine l'esponente in base 2:
$$
    131_{10} = 100000011_2
$$

Si conclude assemblando le 3 parti ricavate:
$$
    +23.75_{10} = 01000001101111100000000000000000_2
$$

\textbf{Nota: } L'esempio proposto segue lo standard IEEE 754 a singola precisone. Lo standard è caratterizzato da:
\begin{itemize}
    \item Singola precisione, $E = 8$ bit e $M = 23$ bit
    \item Doppia precisione, $E= 11$ bit e $M=52$ bit
\end{itemize}


\textbf{Esempio:}  
$-13.25_{10}$:
\[
13.25 = 1101.01_2 = 1.10101_2 \cdot 2^3
\]
\[
S = 1, \quad E = 3+127 = 130 = 10000010_2
\]
\[
M = 10101000000000000000000
\]

Risultato (32 bit):
\[
1 \; 10000010 \; 10101000000000000000000
\]

\clearpage

\section{Logica Booleana}
Un'espressione booleana descrive una condizione, ovvero un'istruzione che, quando valutata, risulta sempre in un valore pari a zero ($0$, falso) o uno ($1$, vero), determinata dagli operatori e dal valore delle costanti o variabili in essa contenute. Per comporre espressioni booleane complesse si utilizzano i seguenti operatori:

\begin{description}
\item[Operatori relazionali] valutano relazioni binarie tra i due operandi:
  \begin{itemize}
  \item [\texttt{<}] minore di
  \item [\texttt{<=}] minore di o uguale uguale a
  \item [\texttt{>}] maggiore di
  \item [\texttt{>=}] maggiore di o uguale a
  \item [\texttt{==}] uguale a
  \item [\texttt{!=}] non uguale a (diverso)
  \end{itemize}
\item[Operatori booleani] valutano condizioni di verit\`a tra i due operandi
  \begin{itemize}
  \item [\texttt{\&\&}] AND (congiunzione logica)
  \item [\texttt{||}] OR (disgiunzione logica)
  \item [\texttt{!}] NOT (negazione logica)
  \end{itemize}
    N.B. Differentemente dalle operazioni AND e OR il NOT \`e un operatore unario, valutato quindi su singolo input.
\end{description}

Le espressioni logiche rispettano un ordine tra i propri operatori. In assenza di parentesi l'ordine è determinato come (elementi a priorità maggiore in alto):
\begin{enumerate}
    \item negazione (NOT) \texttt{!}
    \item operatori di relazione \texttt{<}, \texttt{>}, \texttt{<=}, \texttt{>=}
    \item uguaglianza \texttt{==}, disuguaglianza \texttt{!=},
    \item congiunzione (AND) \texttt{\&\&}
    \item disgiunzione (OR) \texttt{||}
\end{enumerate}

% \paragraph{Attenzione:} si osservi che in C (come in molti altri linguaggi), l'operazione di assegnamento \texttt{a = 3} \`e diversa dall'operazione di confronto \texttt{a == 3}. La prima \`e sempre valutata come vera ($1$, uno), mentre la seconda, ovviamente, dipende dal valore memorizzato in \texttt{a}. Perci\`o:
% \lstinputlisting[style=cstyle]{src/assegnamento_vs_confronto.c}


\textbf{Proprietà}: gli operatori \texttt{\&\&} (AND) e \texttt{||} (OR) sono commutativi:
\begin{itemize}
    \item a \texttt{\&\&} b == b \texttt{\&\&} a
    \item a \texttt{||} b == b \texttt{||} a
\end{itemize}
Le doppie negazioni si elidono:
\begin{itemize}
    \item \texttt{!!}a == a
\end{itemize}

\subsection{Tabelle di Verità}
Le tabelle di verit\`a sono uno strumento logico utilizzato per determinare la validit\`a di espressioni o proposizioni logiche. Esse mostrano tutte le possibili combinazioni di valori di verit\`a (vero o falso) delle variabili coinvolte, permettendo di valutare l'output di un'espressione logica in ciascun caso. La tabella \`e strutturata come segue (esempio in Tabella \ref{tab:logica}):
\begin{itemize}
    \item Una riga per ogni possibile assegnamento di valori logici alle variabili
    \item Una colonna per ogni espressione che compone l\textquotesingle espressione data (inclusa la formula stessa)
\end{itemize}
\begin{table}[h!]
\centering
    \begin{tabular}{|c|c|c|}
        \hline
        \multicolumn{3}{|c|}{\textbf{\large Congiunzione (AND)}} \\
        \hline
        \textbf{A} & \textbf{B} & \textbf{A \&\& B} \\
        \hline
        0 & 0 & 0 \\
        \hline
        0 & 1 & 0 \\
        \hline
        1 & 0 & 0 \\
        \hline
        \rowcolor{yellow!30} 1 & 1 & 1 \\
        \hline
    \end{tabular}
    \caption{Tabella di verit\`a per l'espressione logica A \texttt{\&\&} B}
    \label{tab:logica}
\end{table}

\paragraph{Leggi di De Morgan:} illustrano come distribuire la negazione rispetto a \texttt{||} e \texttt{\&\&}:
\begin{itemize}
    \item \texttt{!}(a \texttt{\&\&} b) == !a \texttt{||} !b
    \item \texttt{!}(a \texttt{||} b) == !a \texttt{\&\&} !b
\end{itemize}

\textbf{Esempio: }
\begin{center}
\begin{tabular}{@{}l@{}}
\texttt{!((a >= 5) \&\& (a <= 10))} \\
\texttt{!(a >= 5) || !(a <= 10)} \quad [De Morgan] \\
\texttt{!!(a < 5) || !!(a > 10)} \quad [trasformazione disuguaglianze] \\
\texttt{(a < 5) || (a > 10)} \quad [doppia negazione]
\end{tabular}
\end{center}
Utilizzando De Morgan rendiamo l'espressione più leggibile. Avremo quindi in questo caso risultato vero se il valore di a è minore di 5 oppure maggiore di 10.
% \subsection{Esercizi}

% \begin{Exercise}[label={ex:tab_1}]
    Si consideri la seguente espressione booleana:
    $$(A \ \&\& \ B) \ || \ !C$$
    e se ne compili la corrispondete tabella di verit\`a.
\end{Exercise}

\begin{Answer}[ref={ex:tab_1}]
\centering
    \begin{tabular}{|c|c|c|c|c|c|}
        \hline
        A & B & C & A \&\& B & !C & (A \&\& B) || !C \\
        \hline
        0 & 0 & 0 & 0 & 1 & 1 \\
        \hline
        0 & 0 & 1 & 0 & 0 & 0 \\
        \hline
        0 & 1 & 0 & 0 & 1 & 1 \\
        \hline
        0 & 1 & 1 & 0 & 0 & 0 \\
        \hline
        1 & 0 & 0 & 0 & 1 & 1 \\
        \hline
        1 & 0 & 1 & 0 & 0 & 0 \\
        \hline
        1 & 1 & 0 & 1 & 1 & 1 \\
        \hline
        1 & 1 & 1 & 1 & 0 & 1 \\
        \hline
    \end{tabular}
\end{Answer} 
% \begin{Exercise}[label={ex:implicazione}]
		\Question Scrivere la tabella di verit\'a della proposizione $A \Rightarrow B$ (implicazione logica), sapendo che essa \`e equivalente alla proposizione $!A \; || \; B$ (due proposizioni sono equivalenti se la loro tabella di verit\`a \`e identica).

		\Question Scrivere la tabella di verit\`a della proposizione $A \Leftrightarrow B$ (co-implicazione logica), sapendo che essa \`e equivalente alla proposizione $A \Rightarrow B \; \&\& \; B \Rightarrow A$.
		
		\Question Mostrare che la proposizione $A \Rightarrow B$ non \`e equivalente alla proposizione $A \; || \; !B$.
\end{Exercise}

\begin{Answer}[ref={ex:implicazione}]

\begin{enumerate}
	\item Tabella di verit\'a di $A \Rightarrow B$:
	\begin{table}[H]
		\centering
			\begin{tabular}{|c | c | c | c|}\hline
				A & B & !A & !A || B \\ \hline
				0 & 0 & 1  & 1       \\ \hline
				0 & 1 & 1  & 1       \\ \hline
				1 & 0 & 0  & 0       \\ \hline
				1 & 1 & 0  & 1      \\ \hline
			\end{tabular}
	\end{table}
	\item Tabella di verit\'a di $A \Leftrightarrow B$:
	\begin{table}[H]
		\centering
			\begin{tabular}{|c | c | c | c | c|}\hline
				A & B & A $\Rightarrow$ B & B $\Rightarrow$ A & A $\Rightarrow$ B \&\& B $\Rightarrow$ A \\ \hline
				0 & 0 & 1  & 1   & 1    \\ \hline
				0 & 1 & 1  & 0  & 0     \\ \hline
				1 & 0 & 0  & 1  & 0     \\ \hline
				1 & 1 & 1 & 1     & 1   \\ \hline
			\end{tabular}
	\end{table}
	\item Confrontiamo le due tabelle di verit\'a:
	\begin{table}[H]
		\centering
			\begin{tabular}{|c | c | c | c | c|}\hline
				A & B & !B & A || !B & A $\Rightarrow$ B \\ \hline
				0 & 0 & 1  & 1  & 1     \\ \hline
				0 & 1 & 0 & 0  & 1   \\ \hline
				1 & 0 & 1  & 1   & 0    \\ \hline
				1 & 1 & 0  & 1     & 1 \\ \hline
			\end{tabular}
	\end{table}
	Le due tabelle di verit\'a non sono identiche, per cui le proposizioni non sono equivalenti.
	
\end{enumerate}

\end{Answer}
% \begin{Exercise}[label={ex:vacanza}]

(TdE Gennaio 2019) Non vado in vacanza se almeno due dei tre seguenti predicati non sono soddisfatti: (A) ho un gruzzoletto da parte, (B) ho almeno due settimane in cui posso non studiare, (C) ho finito gli esami a luglio.

		\Question Scrivere la tabella di verit\`a per il predicato ``Non vado in vacanza''  attraverso i predicati semplici A, B e C.
		
		\Question La traduzione di tale condizione con $((!A \; || \; !B) \; \&\& \; (!A \; || \; !C) \; \&\& \; (!B \; || \; !C))$ \`e corretta?
\end{Exercise}

\begin{Answer}[ref={ex:vacanza}]

\begin{enumerate}
	\item Due dei tre predicati devono essere falsi perch\'e io non vada in vacanza, quindi la traduzione della precedente condizione \'e:
	
			(! A \&\& ! B) || (! A \&\& ! C) || (! B \&\& ! C)

			in tabella della verit\`a:
			
			\begin{table}[H]
\centering
\begin{tabular}{c|c|c|c}
A & B & C & Non vado in vacanza \\ \hline
0 & 0 & 0 & 1                   \\
0 & 0 & 1 & 1                   \\
0 & 1 & 0 & 1                   \\
0 & 1 & 1 & 0                   \\
1 & 0 & 0 & 1                   \\
1 & 0 & 1 & 0                   \\
1 & 1 & 0 & 0                   \\
1 & 1 & 1 & 0                  
\end{tabular}
\end{table}

	\item Per la verifica che il predicato \'e effettivamente corretto possiamo utilizzare le leggi di de Morgan per verificare l'equivalenza con la soluzione del punto 1. In alternativa si pu\`o scrivere la tabella di verit\`a e verificare che rifletta la condizione richiesta.

\end{enumerate}

\end{Answer} %riscrivere
% \begin{Exercise}[label={ex:logica}]
  	Assumendo che \texttt{c1} e \texttt{c2} siano due variabili di tipo char, che memorizzano rispettivamente i valori \verb|'e'| ed \verb|'m'| indicare, per ognuna delle espressioni logiche:
  
	  \begin{enumerate}
			  \item se l'espressione \`e vera o falsa (per i valori delle variabili sopra indicati);
			  \item se \`e sempre vera per qualsiasi valore che le due variabili possono assumere;
			  \item se \`e sempre falsa per qualsiasi valore che le due variabili possono assumere.
	  \end{enumerate}

  	Si fornisca una giustificazione per ogni risposta inserita nella tabella (risposte senza giustificazione potranno essere considerate nulle).

		\begin{enumerate}
			  \item \verb#((c1 != 'e') && (c2 == 'm')) || ((c1 != 'h') && (c2 == 'm'))#
			  \item \verb#(c1 < 'g') || !((c1 <= 'g') && (c1 != 'g'))#
			  \item \verb#(c1 <= 'm') || ((c2 > 'm') || !(c2 > c1))#
		\end{enumerate}
\end{Exercise}

\begin{Answer}[ref={ex:logica}]
\begin{enumerate}
	  \item \`E vera perch\'e \`e vero il secondo disgiunto (di cui \`e vero sia il primo congiunto, in quanto \texttt{c1} vale \verb|'e'|, che il secondo congiunto, in quanto \texttt{c2} vale \verb|'m'|). L'espressione, essendo vera, ovviamente non \`e sempre falsa. Inoltre non \`e sempre vera, per esempio sarebbe falsa se \texttt{c2} valesse \verb|'k'| o qualsiasi altro valore diverso da \verb|'m'|.
  	\item \`E sempre vera perch\'e, in base alla legge di De Morgan, \`e equivalente a
	  \begin{center}
	  		\verb#(c1 <'g') || (c1 > 'g') || (c1 == 'g')#
	  \end{center}
	  quindi la formula non \`e sempre falsa.
	  \item \`E sempre vera perch\'e equivalente alla formula
	    \begin{center}
	  		\verb#(c1 <= 'm') || (c2 > 'm') || (c2 <= c1)#
	  \end{center}
	  che \`e identicamente vera: infatti sarebbe falsa solo se fossero falsi tutti e tre i suoi disgiunti, cosa che non pu\`o essere, perch\'e se sono falsi i primi due (cio\`e se \verb#c1 > 'm'# e \verb#c2 <= 'm'#) allora \verb#c2 < c1#, quindi il terzo \`e vero. Quindi la formula non \`e sempre falsa.
\end{enumerate}
  
\end{Answer} 
% \begin{Exercise}[label={ex:espr_logica}]
Si consideri la seguente espressione booleana:
$$((!A \ || \ B) \ \&\& \ (A \  \&\& !B))\ || \ C$$
e se ne compili la corrispondente tabella di verit\`a (ossia una tabella in cui $0$ rappresenta il valore logico FALSO, $1$ il valore VERO).

Si consideri ora la condizione, scritta in linguaggio C, in cui x e y siano due variabili intere:
$$((!(x < 0) || (y>0)) \ \&\& \ ((x < 0) \ \&\& \ !(y > 0))) ) \ || \ (x > 3)$$
ottenuta dall'espressione booleana sopra indicata sostituendo la variabile $A$ con $x < 0$, la variabile $B$ con $y > 0$, la variabile $C$ con $x > 3$
Si risponda alle seguenti domande giustificando la risposta:
\begin{enumerate}
	\item L'espressione \`e vera o falsa quando $x = 5$ e $y = 7$?
	\item Se $x < 0$, per quali valori di y l'espressione \`e vera?
\end{enumerate}
\end{Exercise}

\begin{Answer}[ref={ex:espr_logica}]
\begin{tabular}{|c|c|c|c|c|c|c|c|c|}
\hline
A & B & C & !A & !B & (!A \ || \ B) & (A \  \&\& !B)) & ((!A \ || \ B) \ \&\& \ (A \  \&\& !B)) & TOT\\
\hline
0 & 0 & 0 & 1 & 1 & 1 & 0 & 0 & 0\\
\hline
0 & 0 & 1 & 1 & 1 & 1 & 0 & 0 & 1\\
\hline
0 & 1 & 0 & 1 & 0 & 1 & 0 & 0 & 0\\
\hline
0 & 1 & 1 & 1 & 0 & 1 & 0 & 0 & 1\\
\hline
1 & 0 & 0 & 0 & 1 & 0 & 1 & 0 & 0\\
\hline
1 & 0 & 1 & 0 & 1 & 0 & 1 & 0 & 1\\
\hline
1 & 1 & 0 & 0 & 0 & 1 & 0 & 0 & 0\\
\hline
1 & 1 & 1 & 0 & 0 & 1 & 0 & 0 & 1\\
\hline
\end{tabular}

Quando $x = 5$, la variabile $C = 1$, quindi l'espressione \`e vera. Se $x < 0$, la variabile $C = 0$, quindi l'espressione \`e falsa per qualsiasi valore di $y$.
\end{Answer}
% \begin{Exercise}[label={ex:cinema}]

(TdE Luglio 2018) Si considerino le seguenti frasi, esprimibili in logica proposizionale tramite tre predicati (A, B, C), che esprimono due condizioni per cui si pu\`o entrare al cinema gratis:

\begin{enumerate}
	\item Sei il millesimo cliente E sei il cliente pi\`u paffutello della sala, OPPURE non hai mai visto il film in programma.
	\item Hai gi\`a visto il film in programma E almeno una delle seguenti condizioni \`e vera:
\begin{itemize}
	\item non sei il millesimo cliente;
	\item c'\`e qualche cliente in sala paffutello almeno quanto te.
\end{itemize}
\end{enumerate}

Si definiscano i tre predicati logici A, B, C contenuti nelle due precedenti frasi, si traducano quindi le due frasi in espressioni booleane, si dica se le proposizioni espresse nelle due frasi sono equivalenti, giustificando la risposta.		
		
\end{Exercise}

\begin{Answer}[ref={ex:cinema}]

	Definiamo:
	\begin{itemize}
		\item A:  ``sei il millesimo cliente'';
		\item B: ``sei il cliente pi\`u paffutello della sala'';
		\item C: ``hai visto il film in programma''.
	\end{itemize}
	
	La prima frase corrisponde a: (A \&\& B) || !C.
	
	La seconda frase corrisponde a: C \&\& (!A || !B).
	
	Le due espressioni non sono equivalenti perch\'e, quando C \'e falsa, la prima \'e sempre vera, mentre la seconda \'e sempre falsa. In alternativa si scrivano le tabelle di verit\`a delle due espressioni e si noti che non sono identiche.

\end{Answer} %riscrivere

\begin{Exercise}[label={ex:conv210}]
		Convertire in base 10 la sequenza $101101_2$ 

\end{Exercise}

\begin{Answer}[ref={ex:conv210}]
    \begin{align*}
      101101_2 &= 1 \cdot 2^5 + 0 \cdot 2^4 + 1 \cdot 2^3 + 1 \cdot 2^2 + 0 \cdot 2^1 + 1 \cdot 2^0 \\
      &= 32 + 0 + 8 + 4 + 0 + 1 \\
      &= 45
    \end{align*}
\end{Answer}
\begin{Exercise}[label={ex:conv210}]
		Convertire in base 2 il valore $543_{10}$
\end{Exercise}

\begin{Answer}[ref={ex:conv210}]
    Per convertire il valore $543_{10}$ saranno necessari $n \geq \lceil \log_2(543) \rceil = \lceil 9.08 \rceil = 10$ bit. 
    \begin{center}
    \begin{tabular}{r|r}
        \multicolumn{2}{c}{2} \\
        543 & 1 \\
        271 & 1 \\
        135 & 1 \\
        67 & 1 \\
        33 & 1 \\
        16 & 0 \\
        8 & 0 \\
        4 & 0 \\
        2 & 0 \\
        1 & 1 \\
        0
    \end{tabular}
    \end{center}
    Il risultato finale è $543_{10} = 1000011111_2$
\end{Answer}
\begin{Exercise}[label={ex:bin_ottale}]
    Convertire il numero binario $(10111001101)_{2}$ in base 8 (ottale).
\end{Exercise}

\begin{Answer}[ref={ex:bin_ottale}]
    Si raggruppano i bit a blocchi di tre, partendo da destra. Si aggiungono zeri a sinistra se necessario (in questo caso, un solo zero):
    \[
    (\underbrace{010}_{2}\ \underbrace{111}_{7}\ \underbrace{001}_{1}\ \underbrace{101}_{5})_{2}
    \]
    Si converte ciascun gruppo di tre bit nel suo equivalente decimale (che è anche la cifra ottale):
    \begin{align*}
        010_2 &= 2 \\
        111_2 &= 4 + 2 + 1 = 7 \\
        001_2 &= 1 \\
        101_2 &= 4 + 1 = 5
    \end{align*}
    Il risultato si ottiene assemblando le cifre ottali:
    \[
    (10111001101)_{2} = (2715)_{8}
    \]
\end{Answer}
\begin{Exercise}[label={ex:bin_esa}]
    Convertire il numero binario $(11011001011)_{2}$ in base 16 (esadecimale).
\end{Exercise}

\begin{Answer}[ref={ex:bin_esa}]
    Si raggruppano i bit in blocchi di 4, partendo da destra, aggiungendo zeri a sinistra per completare l'ultimo blocco:
    \begin{align*}
        (11011001011)_{2} &\to (\underbrace{0110}_{6}\ \underbrace{1100}_{C}\ \underbrace{1011}_{B})_{2} \\
        (11011001011)_{2} &= (6CB)_{16}
    \end{align*}
\end{Answer}

\begin{Exercise}[label={ex:somma_overflow}]
    Eseguire la somma binaria $A = (11011000)_{2}$ e $B = (01101101)_{2}$. Assumendo una rappresentazione a 8 bit per numeri naturali, verificare se si verifica un overflow.
\end{Exercise}

\begin{Answer}[ref={ex:somma_overflow}]
    Si esegue la somma binaria:
    \begin{table}[H]
    \centering
    \begin{tabular}{ccccccccccc}
            &      & 1 & 1 & 1 & 1 & 1 & 0 & 0 & 0 & \\
            &      & 1 & 1 & 0 & 1 & 1 & 0 & 0 & 0 &+\\
            &      & 0 & 1 & 1 & 0 & 1 & 1 & 0 & 1 &=\\
            \hline
            &(1)& 0 & 1 & 0 & 0 & 0 & 1 & 0 & 1 &\\
    \end{tabular}
    \end{table}
    Il risultato è $S = (101000101)_{2}$.

    Poiché la rappresentazione è a 8 bit e il bit di riporto finale ($c_9$) è $\mathbf{1}$, si verifica un \textbf{overflow}.
\end{Answer}

\begin{Exercise}[label={ex:somma_conv}]
    Convertire in base 2 i numeri 77 e 156. Sommare in base 2 i risultati ottenuti.
\end{Exercise}

\begin{Answer}[ref={ex:somma_conv}]
    Per convertire il valore $77_{10}$ saranno necessari $n \geq \lceil \log_2(77) \rceil = \lceil 6.26 \rceil = 7$ bit. 

    \begin{center}
    \begin{tabular}{r|r}
        \multicolumn{2}{c}{2} \\
        77 & 1 \\
        38 & 0 \\
        19 & 1 \\
        9 & 1 \\
        4 & 0 \\
        2 & 0 \\
        1 & 1 \\
        0 &
    \end{tabular}
    \end{center}

    $77_{10} = 1001101_2$
    
    Per convertire il valore $156_{10}$ saranno necessari $n \geq \lceil \log_2(156) \rceil = \lceil 7.28 \rceil = 8$ bit. 

    \begin{center}
    \begin{tabular}{r|r}
        \multicolumn{2}{c}{2} \\
        156 & 0 \\
        78 & 0 \\
        39 & 1 \\
        19 & 1 \\
        9 & 1 \\
        4 & 0 \\
        2 & 0 \\
        1 & 1 \\
        0 &
    \end{tabular}
    \end{center}

    $156_{10} = 10011100_2$
    
    Eseguiamo ora la somma binaria:
    \begin{table}[H]
    \centering
    \begin{tabular}{ccccccccccc}
            &      & 0 & 0 & 1 & 1 & 1 & 1 & 0 & 0 & \\
            &      & 0 & 1 & 0 & 0 & 1 & 1 & 0 & 1 &+\\
            &      & 1 & 0 & 0 & 1 & 1 & 1 & 0 & 0 &=\\
            \hline
            & & 1 & 1 & 1 & 0 & 1 & 0 & 0 & 1 &\\
    \end{tabular}
    \end{table}
    Il risultato è $S = (11101001)_{2}$.
\end{Answer}

\begin{Exercise}[label={ex:sottrazione_underflow_final_v3}]
    Eseguire la sottrazione binaria $A - B$, dove $A = (00110101)_2$ e $B = (01110000)_2$.
    Assumendo una rappresentazione a 8 bit, determinare il risultato e spiegare se si verifica una condizione di \textbf{underflow}.
\end{Exercise}

\begin{Answer}[ref={ex:sottrazione_underflow_final_v3}]
    
    \begin{table}[H]
    \centering
    \begin{tabular}{ccccccccccc}
        & & 1 & 1 & 0 & 0 & 0 & 0 & 0 & \\ % Riga dei prestiti (B_in)
        & & 0 & 0 & 1 & 1 & 0 & 1 & 0 & 1 & $-$ \\ % Minuendo A
        & & 0 & 1 & 1 & 1 & 0 & 0 & 0 & 0 & = \\ % Sottraendo B
        \hline
        & & 1 & 1 & 0 & 0 & 0 & 1 & 0 & 1 & \\ % Risultato S
    \end{tabular}
    \end{table}

    Il settimo bit necessita di un prestito per completare l'operazione. Il prestito verrà richiesto all'ottavo che essendo 0 non può soddisfare, quindi bisognerà "scalare" ancora più a sinistra, ma non abbiamo più alcun bit disponibile. Questo significa che il risultato non può essere rappresentato come un numero positivo a 8 bit: siamo andati in underflow.
    
    Il risultato binario $(11001001)_2$ è un valore errato in questo contesto.
\end{Answer}


\begin{Exercise}[label={ex:ms}]
    Utilizzando 8 bit, rappresentare il numero decimale $-105_{10}$ in Modulo e Segno (MS).
\end{Exercise}

\begin{Answer}[ref={ex:ms}]
    Il numero è negativo, quindi il bit di segno sarà 1.
    Calcoliamo con i restanti 7 bit il valore assoluto, $|-105| = 105_{10}$:
    
    \begin{center}
    \begin{tabular}{r|r}
        \multicolumn{2}{c}{2} \\
        105 & 1 \\
        52 & 0 \\
        26 & 0 \\
        13 & 1 \\
        6 & 0 \\
        3 & 1 \\
        1 & 1 \\
        0 &
    \end{tabular}
    \end{center}

    $105_{10} = 1101001_2$
    
    Unendo il bit di segno al valore assoluto convertito in base due ottienamo come risultato finale $-105_{10} = 11101001_{2\text{MS}}$
\end{Answer}

\begin{Exercise}[label={ex:complemento}]
		\begin{itemize}
				\item[(a)] Si dica qual \`e l'intervallo di valori interi rappresentabile con la codifica in complemento a due a 9 bit.
				\item[(b)] Con riferimento a tale codifica indicare, giustificando brevemente le risposte, quali delle seguenti operazioni possono essere effettuate correttamente:
				\begin{itemize}
						\item $-254 - 255$
						\item $+254 - 253$
						\item $-18 + 236$
						\item $+217 + 182$
				\end{itemize}
				\item[(c)] Mostrare in dettaglio come avviene il calcolo delle operazioni (i) e (ii), evidenziando il bit di riporto e il bit di overflow cos\`i ottenuti.
		\end{itemize}

\end{Exercise}

\begin{Answer}[ref={ex:complemento}]
		\begin{itemize}
				\item[(a)] I valori rappresentabili vanno da $- 2^{m-1} = -256$ a $2^{m-1}-1 = +255$ compresi.
				\item[(b)] Le soluzioni delle operazioni sono:
				\begin{itemize}
						\item $-254 - 255$ NO si ottiene un valore negativo troppo grande in valore assoluto
						\item $+254 - 253$ SI si ottiene un valore piccolo in valore assoluto
						\item $- 18 + 236$ SI si ottiene un valore positivo, grande in valore assoluto ma nei limiti
						\item $+217 + 182$ NO si ottiene un valore positivo troppo grande in valore assoluto
						\end{itemize}
				\item[(c)]
						\item In complemento a 2 a $m = 9$ bit, $(-254)_{10} = (100000010)_{CP2}$ (perch\`e $(254)_{10} = (011111110)_2$). Da questo possiamo ricavare che $(-255)_{10} = (100000001)_{CP2}$ essendo $-255 = -254 - 1$. La somma diventa:\\
						\begin{table}[H]
						\centering
						\begin{tabular}{cccccccccccc}
								&& (1)\\ 
								&& 1 & 0 & 0 & 0 & 0 & 0 & 0 & 1 &0 &+\\
								&& 1 & 0 & 0 & 0 & 0 & 0 & 0 & 0 &1 &=\\
								\hline
								[1]&(1)& 0 & 0 & 0 & 0 & 0 & 0 & 0 & 1 &1 &\\
						\end{tabular}
						\end{table}
						\item $254$ in complemento a due \`e $(011111110)_{CP2} = (011111110)_2$ essendo positivo. Invece $-253 = -254 +1 = (100000011)_{CP2}$, quindi:\\
						\begin{table}[H]
						\centering
						\begin{tabular}{cccccccccccc}
								&      &(1)&(1)&(1)&(1)&(1)&(1)&(1)&(1)&  & \\
								&      & 0 & 1 & 1 & 1 & 1 & 1 & 1 & 1 &0 &+\\
								&      & 1 & 0 & 0 & 0 & 0 & 0 & 0 & 1 &1 &=\\
								\hline
								[0]&(1)& 0 & 0 & 0 & 0 & 0 & 0 & 0 & 0 &1 &\\
						\end{tabular}
						\end{table}
						ignorando il riporto abbiamo il risultato esatto $(000000001)_{CP2} = (1)_{10}$.
				\end{itemize}
\end{Answer}

\begin{Exercise}[label={ex:virgola_fissa}]
    Convertire il numero decimale $+10.625_{10}$ in binario a virgola fissa, utilizzando 5 bit per la parte intera (incluso il segno) e 4 bit per la parte frazionaria.
\end{Exercise}

\begin{Answer}[ref={ex:virgola_fissa}]
    \begin{enumerate}
        \item Parte Intera ($10_{10}$ su 5 bit):
        $$10_{10} = 1010_2$$
        Adattato a 5 bit (segno $0$): $01010$.
        \item Parte Frazionaria ($0.625_{10}$ su 4 bit):
        Si usa il metodo delle moltiplicazioni successive:
        \begin{align*}
            0.625 \times 2 &= 1.25 \quad \to 1 \\
            0.25 \times 2 &= 0.50 \quad \to 0 \\
            0.50 \times 2 &= 1.00 \quad \to 1
        \end{align*}
        Parte frazionaria: $.101_2$. Adattato a 4 bit: $1010$.
        \item Unione delle due parti:
        $+10.625_{10} = 01010.1010_2$
    \end{enumerate}
\end{Answer}

\begin{Exercise}[label={ex:virgola_mobile}]
    Convertire il numero decimale $-0.125_{10}$ nel formato a virgola mobile IEEE 754 a Singola Precisione (32 bit).
\end{Exercise}

\begin{Answer}[ref={ex:virgola_mobile}]
    \begin{enumerate}
        \item Segno $S$: Negativo, $S=1$
        \item Conversione Binaria e Normalizzazione ($+0.125$):
        \begin{align*}
            0.125_{10} &= 0.001_2 \\
            \text{Normalizzazione:} & \quad 1.0 \times 2^{-3} \\
            E_{eff} &= -3.
        \end{align*}
        \item Esponente Traslato (E su 8 bit): Bias = 127.
        \begin{align*}
            E &= E_{eff} + 127 = -3 + 127 = 124_{10} \\
            E &= 01111100_2
        \end{align*}
        \item Mantissa ($M$) su 23 bit: Dalla forma $1.0$, la parte frazionaria è $0$.
        \[ M = 00000000000000000000000 \]
        \item Uniamo $S$ | $E$ | $M$:
        $-0.125_{10} = 1 \quad 01111100 \quad 00000000000000000000000$
    \end{enumerate}
\end{Answer}


\begin{Exercise}[label={ex:tab_1}]
    Si consideri la seguente espressione booleana:
    $$(A \ \&\& \ B) \ || \ !C$$
    e se ne compili la corrispondete tabella di verit\`a.
\end{Exercise}

\begin{Answer}[ref={ex:tab_1}]
\centering
    \begin{tabular}{|c|c|c|c|c|c|}
        \hline
        A & B & C & A \&\& B & !C & (A \&\& B) || !C \\
        \hline
        0 & 0 & 0 & 0 & 1 & 1 \\
        \hline
        0 & 0 & 1 & 0 & 0 & 0 \\
        \hline
        0 & 1 & 0 & 0 & 1 & 1 \\
        \hline
        0 & 1 & 1 & 0 & 0 & 0 \\
        \hline
        1 & 0 & 0 & 0 & 1 & 1 \\
        \hline
        1 & 0 & 1 & 0 & 0 & 0 \\
        \hline
        1 & 1 & 0 & 1 & 1 & 1 \\
        \hline
        1 & 1 & 1 & 1 & 0 & 1 \\
        \hline
    \end{tabular}
\end{Answer}
% \input{01-bin-logica/ex/ex_logica}
\begin{Exercise}[label={ex:espr_logica}]
Si consideri la seguente espressione booleana:
$$((!A \ || \ B) \ \&\& \ (A \  \&\& !B))\ || \ C$$
e se ne compili la corrispondente tabella di verit\`a (ossia una tabella in cui $0$ rappresenta il valore logico FALSO, $1$ il valore VERO).

Si consideri ora la condizione, scritta in linguaggio C, in cui x e y siano due variabili intere:
$$((!(x < 0) || (y>0)) \ \&\& \ ((x < 0) \ \&\& \ !(y > 0))) ) \ || \ (x > 3)$$
ottenuta dall'espressione booleana sopra indicata sostituendo la variabile $A$ con $x < 0$, la variabile $B$ con $y > 0$, la variabile $C$ con $x > 3$
Si risponda alle seguenti domande giustificando la risposta:
\begin{enumerate}
	\item L'espressione \`e vera o falsa quando $x = 5$ e $y = 7$?
	\item Se $x < 0$, per quali valori di y l'espressione \`e vera?
\end{enumerate}
\end{Exercise}

\begin{Answer}[ref={ex:espr_logica}]
\begin{tabular}{|c|c|c|c|c|c|c|c|c|}
\hline
A & B & C & !A & !B & (!A \ || \ B) & (A \  \&\& !B)) & ((!A \ || \ B) \ \&\& \ (A \  \&\& !B)) & TOT\\
\hline
0 & 0 & 0 & 1 & 1 & 1 & 0 & 0 & 0\\
\hline
0 & 0 & 1 & 1 & 1 & 1 & 0 & 0 & 1\\
\hline
0 & 1 & 0 & 1 & 0 & 1 & 0 & 0 & 0\\
\hline
0 & 1 & 1 & 1 & 0 & 1 & 0 & 0 & 1\\
\hline
1 & 0 & 0 & 0 & 1 & 0 & 1 & 0 & 0\\
\hline
1 & 0 & 1 & 0 & 1 & 0 & 1 & 0 & 1\\
\hline
1 & 1 & 0 & 0 & 0 & 1 & 0 & 0 & 0\\
\hline
1 & 1 & 1 & 0 & 0 & 1 & 0 & 0 & 1\\
\hline
\end{tabular}

Quando $x = 5$, la variabile $C = 1$, quindi l'espressione \`e vera. Se $x < 0$, la variabile $C = 0$, quindi l'espressione \`e falsa per qualsiasi valore di $y$.
\end{Answer}
\begin{Exercise}[label={ex:implicazione}]
		\Question Scrivere la tabella di verit\'a della proposizione $A \Rightarrow B$ (implicazione logica), sapendo che essa \`e equivalente alla proposizione $!A \; || \; B$ (due proposizioni sono equivalenti se la loro tabella di verit\`a \`e identica).

		\Question Scrivere la tabella di verit\`a della proposizione $A \Leftrightarrow B$ (co-implicazione logica), sapendo che essa \`e equivalente alla proposizione $A \Rightarrow B \; \&\& \; B \Rightarrow A$.
		
		\Question Mostrare che la proposizione $A \Rightarrow B$ non \`e equivalente alla proposizione $A \; || \; !B$.
\end{Exercise}

\begin{Answer}[ref={ex:implicazione}]

\begin{enumerate}
	\item Tabella di verit\'a di $A \Rightarrow B$:
	\begin{table}[H]
		\centering
			\begin{tabular}{|c | c | c | c|}\hline
				A & B & !A & !A || B \\ \hline
				0 & 0 & 1  & 1       \\ \hline
				0 & 1 & 1  & 1       \\ \hline
				1 & 0 & 0  & 0       \\ \hline
				1 & 1 & 0  & 1      \\ \hline
			\end{tabular}
	\end{table}
	\item Tabella di verit\'a di $A \Leftrightarrow B$:
	\begin{table}[H]
		\centering
			\begin{tabular}{|c | c | c | c | c|}\hline
				A & B & A $\Rightarrow$ B & B $\Rightarrow$ A & A $\Rightarrow$ B \&\& B $\Rightarrow$ A \\ \hline
				0 & 0 & 1  & 1   & 1    \\ \hline
				0 & 1 & 1  & 0  & 0     \\ \hline
				1 & 0 & 0  & 1  & 0     \\ \hline
				1 & 1 & 1 & 1     & 1   \\ \hline
			\end{tabular}
	\end{table}
	\item Confrontiamo le due tabelle di verit\'a:
	\begin{table}[H]
		\centering
			\begin{tabular}{|c | c | c | c | c|}\hline
				A & B & !B & A || !B & A $\Rightarrow$ B \\ \hline
				0 & 0 & 1  & 1  & 1     \\ \hline
				0 & 1 & 0 & 0  & 1   \\ \hline
				1 & 0 & 1  & 1   & 0    \\ \hline
				1 & 1 & 0  & 1     & 1 \\ \hline
			\end{tabular}
	\end{table}
	Le due tabelle di verit\'a non sono identiche, per cui le proposizioni non sono equivalenti.
	
\end{enumerate}

\end{Answer}
\begin{Exercise}[label={ex:de_morgan_tabella}]
    Data l'espressione logica $\text{!(A || !B) \&\& C}$:
    \begin{enumerate}
        \item Utilizzare le Leggi di De Morgan per semplificare la negazione della parentesi.
        \item Costruire la tabella di verità completa per l'espressione originale.
    \end{enumerate}
\end{Exercise}

\begin{Answer}[ref={ex:de_morgan_tabella}]
    \begin{enumerate}
        \item Semplificazione (Leggi di De Morgan):
        \begin{align*}
            \text{!(A || !B) \&\& C} &\equiv (\text{!A \&\& !!B) \&\& C} \quad \text{[De Morgan]} \\
            &\equiv !A \&\& B \&\& C \quad \text{[Doppia Negazione]}
        \end{align*}
        \item Tabella di verità:
            \[
            \begin{array}{|c|c|c|c|c|c|}
                \hline
                \text{A} & \text{B} & \text{C} & \text{!B} & \text{A || !B} & \text{!(A || !B) \&\& C} \\ \hline
                0 & 0 & 0 & 1 & 1 & 0 \\ \hline
                0 & 0 & 1 & 1 & 1 & 0 \\ \hline
                0 & 1 & 0 & 0 & 0 & 0 \\ \hline
                0 & 1 & 1 & 0 & 0 & 1 \\ \hline
                1 & 0 & 0 & 1 & 1 & 0 \\ \hline
                1 & 0 & 1 & 1 & 1 & 0 \\ \hline
                1 & 1 & 0 & 0 & 1 & 0 \\ \hline
                1 & 1 & 1 & 0 & 1 & 0 \\ \hline
        
            \end{array}
            \]
    \end{enumerate}
\end{Answer}

\clearpage
\section*{Soluzioni}
\shipoutAnswer

\end{document}